\documentclass[11pt,a4paper]{article}
\usepackage[utf8]{inputenc}
\usepackage{amsmath,amssymb,amsthm}
\usepackage{listings}
\usepackage[margin=2.5cm]{geometry}
\usepackage{xcolor}

\lstset{
  language=C++,
  basicstyle=\ttfamily\small,
  keywordstyle=\color{blue}\bfseries,
  commentstyle=\color{green!50!black},
  stringstyle=\color{red!70!black},
  numbers=left,
  numberstyle=\tiny\color{gray},
  breaklines=true,
  frame=single,
  tabsize=4,
  showstringspaces=false
}

\title{IOI 1990 -- Problem 1: Subsets}
\author{Solution}
\date{}

\begin{document}
\maketitle

\section{Problem Statement}

Given a set of $n$ integers, find all subsets whose elements sum to a given
target value $T$. Output each such subset.

\section{Solution}

We present two approaches: backtracking for small $n$ and meet-in-the-middle
for moderate $n$.

\subsection{Method 1: Backtracking ($n \le 20$)}

Enumerate subsets recursively, maintaining a running sum. At each position,
either include or exclude the current element. When the sum equals $T$,
output the current subset.

To avoid duplicate subsets when elements repeat, sort the array and skip
consecutive equal elements at the same recursion level.

\subsection{Method 2: Meet in the Middle ($n \le 40$)}

Split the array into halves $A$ (first $\lfloor n/2 \rfloor$ elements) and
$B$ (remaining elements). Enumerate all $2^{|A|}$ subset sums from $A$ and
all $2^{|B|}$ subset sums from $B$. For each sum $s$ from $B$, look up
$T - s$ among the sums from $A$.

\section{Complexity Analysis}

\begin{itemize}
  \item \textbf{Backtracking:} $O(2^n)$ worst case, significantly reduced by
    pruning in practice.
  \item \textbf{Meet in the Middle:} $O(2^{n/2} \cdot n)$ time,
    $O(2^{n/2})$ space.
\end{itemize}

\section{Implementation: Backtracking}

\begin{lstlisting}
#include <iostream>
#include <vector>
#include <algorithm>
using namespace std;

int n, target;
vector<int> arr;
vector<int> current;
int count_solutions = 0;

void backtrack(int idx, int sum) {
    if (sum == target && !current.empty()) {
        cout << "{ ";
        for (int x : current) cout << x << " ";
        cout << "}" << endl;
        count_solutions++;
    }
    if (idx == n) return;
    if (sum > target && arr[idx] > 0) return; // pruning for positive values

    for (int i = idx; i < n; i++) {
        // Skip duplicates at the same level
        if (i > idx && arr[i] == arr[i-1]) continue;
        current.push_back(arr[i]);
        backtrack(i + 1, sum + arr[i]);
        current.pop_back();
    }
}

int main() {
    cin >> n >> target;
    arr.resize(n);
    for (int i = 0; i < n; i++)
        cin >> arr[i];

    sort(arr.begin(), arr.end());
    backtrack(0, 0);

    if (count_solutions == 0)
        cout << "No subsets found." << endl;
    return 0;
}
\end{lstlisting}

\section{Implementation: Meet in the Middle}

\begin{lstlisting}
#include <iostream>
#include <vector>
#include <map>
using namespace std;

int main() {
    int n, target;
    cin >> n >> target;
    vector<int> arr(n);
    for (int i = 0; i < n; i++)
        cin >> arr[i];

    int half = n / 2;
    int other = n - half;

    // Generate all subset sums for the first half
    map<int, vector<vector<int>>> leftSums;
    for (int mask = 0; mask < (1 << half); mask++) {
        int s = 0;
        vector<int> subset;
        for (int i = 0; i < half; i++) {
            if (mask & (1 << i)) {
                s += arr[i];
                subset.push_back(arr[i]);
            }
        }
        leftSums[s].push_back(subset);
    }

    // Generate all subset sums for the second half;
    // look up complement in leftSums
    for (int mask = 0; mask < (1 << other); mask++) {
        int s = 0;
        vector<int> subset;
        for (int i = 0; i < other; i++) {
            if (mask & (1 << i)) {
                s += arr[half + i];
                subset.push_back(arr[half + i]);
            }
        }
        int need = target - s;
        auto it = leftSums.find(need);
        if (it != leftSums.end()) {
            for (auto& left : it->second) {
                if (left.empty() && subset.empty()) continue;
                cout << "{ ";
                for (int x : left) cout << x << " ";
                for (int x : subset) cout << x << " ";
                cout << "}" << endl;
            }
        }
    }

    return 0;
}
\end{lstlisting}

\section{Notes}

\begin{itemize}
  \item The backtracking solution is the most natural for the original IOI
    1990 constraints (small $n$). Sorting enables both pruning and duplicate
    avoidance.
  \item The meet-in-the-middle approach extends the practical range to about
    $n = 40$. Its memory usage of $O(2^{n/2})$ is the main limiting factor.
  \item Both solutions output all valid subsets. If only the count is needed,
    the meet-in-the-middle approach can use a simple frequency map instead of
    storing explicit subsets.
\end{itemize}

\end{document}
